\documentclass{beamer}
\usepackage{graphicx}
\usepackage{xcolor}
\usepackage{tikz}

% 自定义颜色
\definecolor{purpleblue}{RGB}{80, 80, 180}
\definecolor{lightgray}{RGB}{230, 230, 230}
\definecolor{novelPurple}{RGB}{98, 54, 255} 
\definecolor{darkgray}{RGB}{50, 50, 50}

% 主题设置
\usetheme{Madrid}  
\usecolortheme{dolphin} 

% 设置字体（移除 fontspec，确保 pdfLaTeX 兼容）
\setbeamerfont{title}{size=\Large,series=\bfseries}
\setbeamerfont{frametitle}{size=\LARGE,series=\bfseries}
\setbeamercolor{frametitle}{fg=white, bg=novelPurple}
\setbeamercolor{title}{fg=white, bg=purpleblue}

% 封面信息
\title[SAT Hard Question]{\textbf{SAT Hard Question 16}}
\author{\textbf{Jack Zeng}}
\institute{\textcolor{white}{\textbf{Novel Prep}}}
\date{\today} 

\begin{document}

% --- Title Slide ---
\begin{frame}
    \titlepage
    \vfill
    \begin{flushright}
        \textcolor{purpleblue}{\Huge \textbf{Novel Prep}}
    \end{flushright}
\end{frame}


\begin{frame}
    \begin{tikzpicture}[remember picture, overlay]
        % 设置浅色渐变背景
        \shade[top color=softPurple, bottom color=softGray] 
            (current page.north west) rectangle (current page.south east);
        
        % 添加高斯模糊光晕
        \fill[white, opacity=0.3] (current page.center) circle (4cm);
        
        % 主标题
        \node[align=center] at (current page.center) {
            {\Huge\bfseries\textcolor{purpleblue}{Novel Prep}} \\[0.5cm]
            {\Large\itshape\textcolor{lightText}{Excellence in SAT Prep}}
        };

        % 添加柔和光点（点缀）
        \foreach \x in {1,2,3,4,5} {
            \fill[purpleblue, opacity=0.2] (rand*10-5, rand*7-3.5) circle (0.1+\x*0.05);
        }

        ;
    \end{tikzpicture}
\end{frame}



\begin{frame}{\textbf{Question: Consecutive Integers and Difference Condition}}
\small
For a set of three consecutive integers, the result of subtracting the greatest integer from the sum of the other two integers is $48$. What is the value of the greatest of the three integers?
\end{frame}

\begin{frame}{\textbf{Answer Explanation: Algebraic Setup and Solution}}
\small
Let the three consecutive integers be $n$, $n+1$, and $n+2$.

Sum of the smaller two integers:
\[
n + (n + 1) = 2n + 1
\]

Subtract the greatest integer:
\[
2n + 1 - (n + 2) = 2n + 1 - n - 2 = n - 1
\]

Set equal to $48$:
\[
n - 1 = 48
\]
\[
n = 49
\]

The greatest integer is:
\[
n + 2 = 49 + 2 = \boxed{51}
\]
\end{frame}

\begin{frame}{\textbf{Question: Quadratic Function and Ground Impact Time}}
\small
The function $h(t) = -16t^2 + b$ estimates an object's height, in feet, above the ground $t$ seconds after the object is dropped, where $b$ is a constant.

The function estimates that the object is $40.96$ feet above the ground when it is dropped at $t = 0$.

How many seconds after being dropped does the function estimate the object will hit the ground?
\end{frame}

\begin{frame}{\textbf{Answer Explanation: Solving the Quadratic Equation}}
\small
Since $t = 0$ gives $h(0) = b = 40.96$, the function is:
\[
h(t) = -16t^2 + 40.96
\]

Set $h(t) = 0$ to find when the object hits the ground:
\[
-16t^2 + 40.96 = 0
\]
\[
16t^2 = 40.96
\]
\[
t^2 = \frac{40.96}{16} = 2.56
\]
\[
t = \sqrt{2.56} = 1.6
\]

The object hits the ground after:
\[
\boxed{1.6} \text{ seconds}
\]
\end{frame}


\begin{frame}{\textbf{Question: Survey Percentages and Totals}}
\small
A researcher surveyed undergraduate students, graduate students, and postdoctoral students. The number of undergraduate students surveyed was $7,150\%$ of the number of postdoctoral students surveyed, and the number of graduate students surveyed was $45\%$ of the number of undergraduate students surveyed.

If there were $6,435$ graduate students surveyed, what was the sum of the number of undergraduate students and postdoctoral students surveyed?
\end{frame}

\begin{frame}{\textbf{Answer Explanation: Setting Up Proportions and Solving}}
\small
Let the number of postdoctoral students be $x$.

Number of undergraduate students:
\[
71.5x
\]

Number of graduate students:
\[
0.45 \times 71.5x = 32.175x
\]

Set equal to $6,435$:
\[
32.175x = 6435
\]
\[
x = \frac{6435}{32.175} = 200
\]

So, number of undergraduate students:
\[
71.5 \times 200 = 14,300
\]

Sum of undergraduate and postdoctoral students:
\[
14,300 + 200 = \boxed{14,500}
\]
\end{frame}

\begin{frame}{\textbf{Question: Parallel Lines Condition for No Solution}}
\small
\[
16x - 20y = 6y + 8
\]
\[
ty = 8x + \frac{1}{5}
\]

In the given system of equations, $t$ is a constant. If the system has no solution, what is the value of $t$?
\end{frame}

\begin{frame}{\textbf{Answer Explanation: Comparing Slopes for No Solution Condition}}
\small
Rewrite both equations in slope-intercept form.

First equation:
\[
16x - 20y = 6y + 8
\]
\[
16x - 20y - 6y = 8
\]
\[
16x - 26y = 8
\]
\[
y = \frac{16}{26}x - \frac{8}{26} = \frac{8}{13}x - \frac{4}{13}
\]

Second equation:
\[
ty = 8x + \frac{1}{5}
\]
\[
y = \frac{8}{t}x + \frac{1}{5t}
\]

For no solution, slopes must be equal:
\[
\frac{8}{13} = \frac{8}{t}
\]
\[
t = 13
\]

So the value of $t$ is:
\[
\boxed{13}
\]
\end{frame}

\begin{frame}{\textbf{Question: Unit Conversion of Bone Mineral Density}}
\small
The bone mineral density of a certain person's thoracic spine was estimated to be $0.812$ grams per square centimeter.

What is this estimated bone mineral density in grams per square millimeter? (1 centimeter = 10 millimeters)
\end{frame}

\begin{frame}{\textbf{Answer Explanation: Converting Area Units}}
\small
\[
1 \text{ cm}^2 = (10 \text{ mm})^2 = 100 \text{ mm}^2
\]

So,
\[
0.812 \text{ g/cm}^2 = \frac{0.812}{100} \text{ g/mm}^2 = 0.00812 \text{ g/mm}^2
\]

Thus, the answer is:
\[
\boxed{0.00812}
\]
\end{frame}


\begin{frame}{\textbf{Question: Product of Roots and Factored Form}}
\small
\[
\frac{5}{9}(5x + 9)(x + \sqrt{5k+9} )(x - \sqrt{5k+9} ) = 0
\]

In the given equation, $k$ is a positive constant. The product of the solutions to the equation is $72$. What is the value of $k$?
\end{frame}

\begin{frame}{\textbf{Answer Explanation}}
\small
For the solution
\[x=-\frac{9}{5} , -\sqrt{5k+9}, and \sqrt{5k+9}\]

The product of the solution is 
\[\frac{9}{5}(5k+9)=72\]

Where \boxed{k=6.2}

\end{frame}

\begin{frame}{\textbf{Question: Surface Area of a Right Rectangular Prism}}
\small
A right rectangular prism has a base area of $56t$ square centimeters. The length of the base is $\frac{14}{3}$ cm, and the height is $3$ cm.

What is the expression of the surface area, in cm$^2$, of the prism?
\end{frame}

\begin{frame}{\textbf{Answer Explanation: Surface Area Calculation}}
\small
\boxed{184t+28}


\end{frame}

\begin{frame}{\textbf{Question: Unit Circle, Coordinates, and Trigonometric Ratio}}
\small
In the $xy$-plane, a unit circle with center at the origin $O$ contains point $A$ at $(1, 0)$ and point $B$ at $\left( \frac{5}{\sqrt{34}}, -\frac{3}{\sqrt{34}} \right)$.

If $\angle AOB$ measures $p$ radians, what is the value of $\frac{\cos p}{\sin p}$?
\end{frame}

\begin{frame}{\textbf{Answer Explanation}}
\small
We are asked to compute:
\[
\frac{\cos p}{\sin p}
\]

Recall the identity:
\[
\frac{\cos p}{\sin p} = \cot p = \frac{1}{\tan p}
\]

Point $B$ has coordinates:
\[
\left( \frac{5}{\sqrt{34}}, -\frac{3}{\sqrt{34}} \right)
\]

So:
\[
\tan p = \frac{\text{y-coordinate}}{\text{x-coordinate}} = \frac{-\frac{3}{\sqrt{34}}}{\frac{5}{\sqrt{34}}} = -\frac{3}{5}
\]

Thus:
\[
\frac{\cos p}{\sin p} = \frac{1}{\tan p} = \frac{1}{ -\frac{3}{5} } = -\frac{5}{3}
\]

The value is:
\[
\boxed{-\frac{5}{3}}
\]
\end{frame}



\begin{frame}{\textbf{Question: Maximum Root of a Quadratic Function}}
\small
The function $h$ is defined by:
\[
h(x) = -\sqrt{x^2+ bx + c} 
\]
where $b$ and $c$ are constants. In the $xy$-plane, the graph of $y = h(x)$ contains the points $(2, 0)$ and $\left( 0, -\sqrt{334} \right)$.

If $h(m) = 0$, what is the greatest possible value of $m$?
\end{frame}

\begin{frame}{\textbf{Answer Explanation: Solving for Maximum Root}}
\small


Thus, the greatest possible value of $m$ is:
\[
\boxed{167}
\]
\end{frame}

\begin{frame}{\textbf{Question: Linear Cost Model for Computer Repair}}
\small
A computer repair specialist charges \$160 for the first two hours of repair plus an hourly fee for each additional hour. The total cost for $6$ hours of repair is \$360.

What is the function $f$ that gives the total cost, in dollars, for $x$ hours of repair, where $x \geq 2$?
\end{frame}

\begin{frame}{\textbf{Answer Explanation: Constructing the Piecewise Linear Model}}
\small
The first two hours cost \$160. For $6$ hours, the total cost is \$360, so the cost of the $4$ additional hours is:
\[
360 - 160 = 200
\]

Hourly fee:
\[
\frac{200}{4} = 50
\]

Model:
\[
f(x) = 160 + 50(x - 2) = 50x + 60
\]

Thus, the function is:
\[
\boxed{f(x) = 50x + 60}
\]
\end{frame}




\end{document}